\section{Spherical extension and Kakeya tubes}\label{sec:kakeya}

Write \(d=n-1\), \(p=\pc=2n/d\), and \(L(R)=\log(2+R)\). We keep the
logarithmic power \(A\) in the wave estimate as a parameter. To see what
happens to it in restriction, we take an incoming spherical wave whose
amplitude varies with direction. Near its focus the radial oscillation
cancels, leaving the extension of that amplitude from the unit sphere.

For \(g\) on \(S^{n-1}\) write
\(Eg(X)=\int_{S^{n-1}}e^{\mathrm iX\cdot\omega}g(\omega)\dd\sigma(\omega)\).
Let \(K_\delta f(\omega)\) be the supremum of the averages of \(|f|\)
over unit tubes of radius \(\delta\) parallel to \(\omega\).

\begin{corollary}[Spherical extension and Kakeya]\label{cor:quant-kakeya}
For \(n\ge2\), \(R\ge2\), \(z\in\RR^n\), and \(0<\delta\le1\),
\begin{align*}
 \|Eg\|_{L^{\pc}(B(z,R))}
      &\le C_n [\log(2+R)]^{2/\pc}\|g\|_{L^{\pc}(S^{n-1})},\\
 \|K_\delta f\|_{L^n(S^{n-1})}
      &\le C_n [\log(2/\delta)]^{4/\pc}\|f\|_{L^n(\RR^n)}.
\end{align*}
Every subset of \(\RR^n\) containing a unit line segment in each direction has Hausdorff dimension \(n\).
\end{corollary}

A wave estimate with logarithmic power \(A\) gives the same power for
spherical extension and power \(2A\) for the Kakeya maximal operator.
The radial example after the proof shows that the latter power must be
at least \(1-1/n\). Thus an unrestricted wave estimate with power
\(1/\pc\) would give the sharp Kakeya maximal logarithm.

In the plane, H\"ormander's estimate gives the sharp logarithmic
power $1/4$ for local circle extension: take the output exponent
$q=4+1/\log R$ in \cite[Corollary~1.3]{Hormander1973FLp} and apply
H\"older's inequality on the observation ball and on the circle.
Constant density forces this power by stationary phase. The planar directional
Kakeya maximal bound has sharp power $1/2$
\cite[Equation~(1.5)]{Bourgain1991Kakeya}; see also C\'ordoba's
theorem for rectangles \cite{Cordoba1977}. Davies proved full
Hausdorff dimension in the plane \cite{Davies1971}, and Wang and
Zahl proved it in three dimensions \cite[Theorem~1.1]{WangZahl2025};
see also \cite{GuthWangZahl2026Kakeya}. A separate three-dimensional
maximal estimate with $\delta^{-\varepsilon}$ loss and a
four-dimensional Hausdorff dimension theorem appear in
\cite[Theorem~1.1]{OpenAI2026KakeyaThree} and
\cite[Theorem~1.1]{OpenAI2026Kakeya4Hausdorff}, respectively.
Here the implication from a wave estimate to a fixed logarithmic
maximal estimate works in every dimension.

\begin{lemma}\label{lem:kakeya-local-extension}
Suppose \eqref{eq:main} holds with \(L(N)^A\) for some \(A\ge0\), for a
smooth radial cutoff \(\chi(|\xi|/N)\) with \(\chi\) supported in
\((1/2,2)\) and equal to one near \(1\). For \(R\ge2\), \(z\in\RR^n\),
and \(g\in L^p(S^{n-1})\), the spherical extension
\(Eg(X)=\int_{S^{n-1}}e^{\mathrm iX\cdot\omega}g(\omega)\dd\sigma(\omega)\)
satisfies
\begin{equation}\label{eq:kakeya-local-extension}
             \|Eg\|_{L^p(B(z,R))}
                   \lesssim L(R)^A\|g\|_{L^p(S^{n-1})}.
\end{equation}
\end{lemma}

\begin{proof}
Choose a nonnegative nonzero \(\psi\in C_c^\infty((1,2))\).
Modulation in time selects a frequency shell of width of order one:
\[
 M_N f=\int \psi(T)e^{\mathrm iNT}U(T)\chi(|D|/N)f\dd T,\qquad
 \|M_N f\|_p\lesssim L(N)^A\|f\|_p
\]
by H\"older in \(T\). For \(1/2\le |Z|\le3\) its convolution kernel is
\begin{equation}\label{eq:kakeya-shell-kernel}
 K_N(Z)=c_n N^{d/2}|Z|^{-d/2}e^{\mathrm iN|Z|}\psi(|Z|)
                         +O(N^{d/2-1}),\qquad c_n\ne0 .
\end{equation}
Indeed, the radial multiplier is \(\chi(\rho/N)\widehat\psi(\rho-N)\).
In polar coordinates the angular integral at \(u=\rho|Z|\) is
\[
 \int_{S^{n-1}}e^{\mathrm iu\omega_1}\dd\sigma(\omega)
   =u^{-d/2}(c_+e^{\mathrm iu}+c_-e^{-\mathrm iu})
                       +O(u^{-d/2-1}),\qquad c_+c_-\ne0 ,
\]
by stationary phase at the two poles. After radial integration the
remainder is \(O(N^{d/2-1})\). In the leading terms one can replace
\(\chi(\rho/N)\rho^{d/2}\) by \(N^{d/2}\) with the same error:
the prefactor, extended by zero, has derivative \(O(N^{d/2-1})\)
and \(\int |s\widehat\psi(s)|\dd s<\infty\). Fourier inversion gives
\(\psi(\pm|Z|)\); only the positive argument contributes.

We linearize the distance phase at the focus, as in H\"ormander's passage
from oscillatory integrals to extension \cite[\S3, (3.8)--(3.9)]{Hormander1973FLp};
compare Tao's deduction of local restriction from Bochner--Riesz estimates
\cite[proof of Theorem~1.1]{Tao1999}.
Take \(f_N(r\omega)=e^{-\mathrm iNr}\psi(r)g(\omega)\), so that
\(\|f_N\|_p=C_\psi\|g\|_p\). For \(|X|\le R\) and \(N\ge4R\),
\[
 N|r\omega-X/N|=Nr-X\cdot\omega+O(R^2/N)
 \qquad (r\in\supp\psi).
\]
Substitute this expansion in \eqref{eq:kakeya-shell-kernel} and
integrate in \(r,\omega\). The amplitude changes by \(O(R/N)\), and
the radial phases cancel:
\[
 N^{-d/2}M_N f_N(X/N)
       =c' Eg(-X)+O\big((1+R^2)N^{-1}\|g\|_1\big),
 \qquad c'=c_n\int r^{d/2}\psi(r)^2\dd r\ne0 .
\]
Choose \(N=(4+R)^{3+n/p}\). The error has \(L^p(B(0,R))\) norm at
most \(C\|g\|_p\). On the left the change of variables costs
\(N^{n/p-d/2}=1\), and \(L(N)\lesssim_n L(R)\). This proves
\eqref{eq:kakeya-local-extension} at \(z=0\). Replacing \(g(\omega)\)
by \(e^{\mathrm iz\cdot\omega}g(\omega)\) gives every center.
\end{proof}

\subsection{Global extension from the sphere and the truncated cone}
\label{sec:strict-extension}

Lemma~\ref{lem:kakeya-local-extension} gives global estimates at strictly
larger output exponents. For this step any fixed power of the logarithm
suffices; the tube argument in the next subsection will retain that power.
Write \(E\) for extension from \(S^{n-1}\).
For the truncated cone
\(\Sigma=\{(\rho\omega,\rho):1\le\rho\le2,\ \omega\in S^{n-1}\}\),
let \(E_\Sigma G(x,t)=\int_\Sigma e^{\mathrm i(x\cdot\xi+t\tau)}
G(\xi,\tau)\,\dd\Sigma(\xi,\tau)\), with induced surface measure.

\begin{corollary}[Global Fourier extension]\label{cor:spherical-restriction}\label{cor:cone-extension}
Let \(1<a\le\infty\) and \(b<\infty\). If
\[
 b>\pc,\qquad b>\frac{n+1}{n-1}a',
\]
then
\begin{align}
 \|Eg\|_{L^b(\RR^n)}
 &\lesssim_{n,a,b} \|g\|_{L^a(S^{n-1})},\label{eq:spherical-extension-mixed}\\
 \|E_\Sigma G\|_{L^b(\RR^{n+1})}
 &\lesssim_{n,a,b} \|G\|_{L^a(\Sigma)}.\label{eq:cone-extension}
\end{align}
These are strong estimates, including when \(a=\infty\).
For \(b=\infty\) both estimates hold for every \(1\le a\le\infty\).
\end{corollary}

The two strict inequalities give the interior of the conjectured range
for these compact surfaces \cite[Conjecture~1.1(i)]{Nicola2009}.
The three-dimensional sphere estimate in
\cite[\S\S9.1--9.2]{OpenAI2026SphereRestriction} gives this range for
\(S^2\) and, by randomizing caps, an \(L^3\) Kakeya maximal bound with
loss \(\delta^{-\varepsilon}\).

\begin{proof}
The local spherical estimate in Lemma~\ref{lem:kakeya-local-extension},
with \(A=2/\pc\) by Theorem~\ref{thm:main},
costs \(O_\varepsilon(R^\varepsilon)\) for every \(\varepsilon>0\).
By duality it bounds restriction of functions supported in \(B(0,R)\)
from \(L^{\pc'}(\RR^n)\) to \(L^{\pc'}(S^{n-1})\) with that loss.
Since \(1<\pc'<2\), Tao's epsilon-removal theorem
\cite[Theorem~1.2]{Tao1999} gives global diagonal restriction at
every \(1<u<\pc'\). Dualizing gives \(E:L^s\to L^s\) for every
\(\pc<s<\infty\).

If \(a\ge b\), take \(s=b\) and use the finite measure of the sphere.
If \(a<b\), put \(s=1+b/a'\); then \(\pc<s<b\), since
\(\pc=1+(n+1)/(n-1)\). Interpolating \(E:L^1\to L^\infty\) with
\(E:L^s\to L^s\) at \(\theta=s/b\) gives \eqref{eq:spherical-extension-mixed}:
\(\theta/s=1/b\) and \(1-\theta+\theta/s=1/a\).

For the cone use Nicola's slicing argument
\cite[Theorem~1.2(a)(i)]{Nicola2009}. Write \(G_\rho(\omega)=G(\rho\omega,\rho)\).
Since \(\dd\Sigma=\sqrt2\,\rho^{n-1}\dd\rho\,\dd\sigma\),
\[
 E_\Sigma G(x,t)=\sqrt2\int_1^2e^{\mathrm it\rho}\rho^{n-1}
                   E(G_\rho)(\rho x)\,\dd\rho .
\]
Hausdorff--Young in \(t\), Minkowski's inequality (\(b\ge b'\)),
and spatial dilation give
\[
 \|E_\Sigma G\|_{L^b_{x,t}}
 \lesssim\left(\int_1^2
       [\rho^{n-1-n/b}\|E(G_\rho)\|_{L^b_x}]^{b'}\,\dd\rho\right)^{1/b'}
 \lesssim\left\|\,\|G_\rho\|_{L^a(S^{n-1})}\,\right\|_{L^{b'}_\rho([1,2])}.
\]
H\"older's inequality gives \eqref{eq:cone-extension} because \(b>a'\)
implies \(a>b'\). For \(b=\infty\) the defining integrals are bounded
by the \(L^1\) norms on the two compact surfaces.
\end{proof}


\subsection{From caps to tubes}

Let \(T_\delta(\omega,x)\) be the open \(\delta\)-neighborhood of
\(x+[-1/2,1/2]\omega\). Its volume is comparable to \(\delta^d\) for
\(0<\delta\le1\). Thus \(K_\delta f(\omega)\) is the supremum over
\(x\) of \(|T_\delta(\omega,x)|^{-1}\int_{T_\delta(\omega,x)}|f|\).
These averages are continuous in the center and direction for
\(f\in L^n\), by continuity of translations and rotations of the tube
indicator in \(L^{n/(n-1)}\); in particular the supremum is measurable.

A cap of radius \(\delta\) produces a wave of roughly constant modulus
on a tube of length \(\delta^{-2}\) and radius \(\delta^{-1}\).
The restriction argument of Beckner, Carbery, Semmes and Soria
\cite{BCSS1989Restriction} and Bourgain \cite{Bourgain1991Kakeya}
puts independent signs on disjoint caps. The expected squared modulus
of the sum is then the sum of the squared moduli. This accounts for
the factor of two in the logarithmic exponent.

\begin{lemma}\label{lem:kakeya-overlap}
Assume \eqref{eq:kakeya-local-extension} and set \(q=p/2=n/(n-1)\).
Let \(0<\delta\le1/2\), let the \(\omega_j\) be \(\delta\)-separated
on \(S^{n-1}\), and choose \(x_j\in\RR^n\) arbitrarily. Then for
\(T_j=T_{2\delta}(\omega_j,x_j)\) and \(a_j\ge0\),
\begin{equation}\label{eq:kakeya-overlap}
 \Big\|\sum_j a_j\mathbf1_{T_j}\Big\|_{L^q(\RR^n)}
  \lesssim L(1/\delta)^{2A}
              \Big(\delta^d\sum_j a_j^q\Big)^{1/q}.
\end{equation}
\end{lemma}

\begin{proof}
Take the disjoint caps \(C_j=\{|\omega-\omega_j|<c\delta\}\) for a
small fixed \(c>0\), and put
\(h_j(\omega)=e^{-\mathrm i\delta^{-2}x_j\cdot\omega}\mathbf1_{C_j}(\omega)\).
For \(Y\in T_j\), write \(Y-x_j=t\omega_j+v\) with \(|t|\le1/2\),
\(|v|<2\delta\). Since
\(\omega_j\cdot(\omega-\omega_j)=-|\omega-\omega_j|^2/2\),
the phase in \(Eh_j(\delta^{-2}Y)\) differs from its value at
\(\omega_j\) by at most \(c^2/4+2c\). Thus
\(|Eh_j(\delta^{-2}Y)|\gtrsim\sigma(C_j)\asymp\delta^d\) on \(T_j\).

Group the centers by the unit lattice cubes \(Q\), with
\(J_Q=\{j:x_j\in Q\}\) (use half-open cubes). The tubes in this group
lie in the concentric cube \(5Q\), so their dilates all lie in a
ball of radius \(C_n\delta^{-2}\). For independent random signs
\(\varepsilon_j\) set \(g_\varepsilon=\sum_{j\in J_Q}
\varepsilon_j a_j^{1/2}h_j\). Disjointness of the caps gives
\(\|g_\varepsilon\|_p^p\lesssim\delta^d\sum_{j\in J_Q}a_j^q\).
At each point Jensen's inequality gives
\(\mathbb E|Eg_\varepsilon|^p\ge
(\sum_{j\in J_Q}a_j|Eh_j|^2)^q\).
Integrate on \(\delta^{-2}(5Q)\), change variables back to \(Y\),
and apply \eqref{eq:kakeya-local-extension} on the enclosing ball:
\[
 \delta^{dp-2n}\int_{5Q}
       \Big(\sum_{j\in J_Q}a_j\mathbf1_{T_j}\Big)^q
 \lesssim L(1/\delta)^{Ap}\,\delta^d\sum_{j\in J_Q}a_j^q .
\]
The prefactor is one because \(dp=2n\). The cubes \(5Q\) have bounded
overlap, so the \(q\)-th power of the norm on the left of
\eqref{eq:kakeya-overlap} is bounded by a constant times the sum of
these integrals. Taking the \(q\)-th root gives \(Ap/q=2A\).
\end{proof}

\subsection{The maximal estimate and Hausdorff dimension}

\begin{proof}[Proof of Corollary~\ref{cor:quant-kakeya}]
Theorem~\ref{thm:main} gives \(A=2/p\) in
Lemma~\ref{lem:kakeya-local-extension}. We keep \(A\) in the next
steps to record the dependence on the wave estimate.
For \(0<\delta\le1/2\) choose a maximal \(\delta\)-separated set
\(\{\omega_j\}\). If \(|\omega-\omega_j|\le\delta\), then
\(T_\delta(\omega,x)\subset T_{2\delta}(\omega_j,x)\). Hence
\[
 \|K_\delta f\|_{L^n(S^{n-1})}
       \lesssim \Big(\delta^d\sum_j K_{2\delta}f(\omega_j)^n\Big)^{1/n}.
\]
To bound this discrete norm, test by \(a_j\ge0\) with
\(\delta^d\sum_j a_j^q=1\). For \emph{any} choice of the centers,
\[
 \delta^d\sum_j\frac{a_j}{|T_j|}\int_{T_j}|f|
 \lesssim \int |f|\sum_j a_j\mathbf1_{T_j}
 \lesssim L(1/\delta)^{2A}\|f\|_n
\]
by Lemma~\ref{lem:kakeya-overlap} and H\"older. Taking suprema over
the centers and then the weights proves the maximal estimate.
For \(\delta\in[1/2,1]\) it follows from H\"older on each tube.

Suppose \(F\) contains a unit segment in every direction. Its open
\(\delta\)-neighborhood \(F_\delta\) contains the entire tube around
each such segment. Consequently
\begin{equation}\label{eq:kakeya-volume}
 |F_\delta|\gtrsim L(1/\delta)^{-2An}
                 =L(1/\delta)^{-2(n-1)}\qquad (A=2/p).
\end{equation}
For bounded sets this already gives full lower and upper Minkowski
dimension. To obtain Hausdorff dimension we have to allow balls of
different sizes in the same cover. In fact the logarithmic rate gives
infinite Hausdorff measure for every gauge
\[
 h_\gamma(r)=r^n[\log(1/r)]^{2An}
                  [\log\log(1/r)]^\gamma,\qquad \gamma>n-1
 \quad (r\ \text{small}).
\]

Since \(h_\gamma(2r)\lesssim h_\gamma(r)\), it suffices to use ball
covers. Consider an open ball cover \(\{B_j\}\) of \(F\) with radii
\(r_j\le e^{-2^m}\). Group the radii by
\(J_k=\{j:2^k\le\log(1/r_j)<2^{k+1}\}\) as in the planar argument of
Keich \cite[Claim 3.2]{Keich1999}. Put
\(\delta_k=e^{-2^{k+1}}\), \(G_k=\bigcup_{j\in J_k}2B_j\) and
\(v_k=\sum_{j\in J_k}r_j^n\). On a segment in direction \(\omega\),
the length covered by group \(k\) is at most
\(C K_{\delta_k}\mathbf1_{G_k}(\omega)\). Indeed, at each covered
point the transverse disk of radius \(\delta_k\) lies in a doubled
ball, and disks at different points give disjoint cross sections of
the tube. Summing lengths therefore gives
\(1\lesssim\sum_{k\ge m}K_{\delta_k}\mathbf1_{G_k}(\omega)\) in
every direction. The maximal estimate and H\"older in \(k\) yield
\[
 1\lesssim\sum_{k\ge m}2^{2Ak}v_k^{1/n}
 \lesssim\Big(\sum_j h_\gamma(r_j)\Big)^{1/n}
              \Big(\sum_{k\ge m} k^{-\gamma/(n-1)}\Big)^{(n-1)/n}.
\]
The last factor tends to zero. Thus the Hausdorff measure for this
gauge is infinite, and \(h_\gamma(r)=o(r^s)\) for every \(s<n\)
gives \(\dim_{\mathrm H}F=n\). In the plane the same calculation with
the classical $L^2$ maximal power $1/2$ cited above gives infinite
Hausdorff measure for
$r^2\log(1/r)[\log\log(1/r)]^\gamma$ whenever $\gamma>1$.
The same covering argument works if a segment in
each of a positive measure set of directions meets \(F\) in length
at least a fixed \(\lambda>0\).
\end{proof}

The radial test \(f_\delta(X)=|X|^{-1}\mathbf1_{\{2\delta<|X|<1/2\}}\)
shows the scale of a possible improvement. Its \(L^n\) norm is
comparable to \([\log(1/\delta)]^{1/n}\) for small \(\delta\), but
the average on a tube centered at the origin in any direction is at
least \(c\int_{4\delta}^{1/4}t^{-1}\dd t\). Hence the maximal
logarithm must have power at least \(1-1/n=2/p\). A wave estimate
with the sharp power \(A=1/p\) for arbitrary data would attain this
bound and give \(|F_\delta|\gtrsim L(1/\delta)^{-(n-1)}\). There are
compact Kakeya sets satisfying the opposite inequality, by
Fern\'andez-Delgado and de la Salle
\cite[Theorem 1.3]{FernandezDelgadoDeLaSalle2026}.

\subsection{Two consequences for other families of segments}

For the Nikodym maximal function \(N_\delta f(x)\), take the supremum
of the same tube averages over all unit segments \emph{through} \(x\).
The estimate just proved gives
\begin{equation}\label{eq:kakeya-nikodym}
 \|N_\delta f\|_{L^n(\RR^n)}
             \lesssim L(1/\delta)^{2A}\|f\|_{L^n(\RR^n)}.
\end{equation}
We give the projective argument to keep track of the logarithm;
compare Tao \cite[Theorem 4.10]{Tao1999} and Gao--Liu--Xi
\cite[\S2.2]{GaoLiuXi2025}. First consider only the
\(\rho\)-neighborhoods of the segments
\(\{x+u\omega:1/2\le u\le1\}\), and denote the supremum of their
averages by \(D_\rho f(x)\). By finitely many rotations we can
restrict to \(|\omega_n|\ge c_n\). Place \(x=(x',t_0)\) in a unit
cube \(Q\) centered at zero and fix \(t_0\). For \(\rho\) small, on each
such tube \(c_n/4\le|s-t_0|\le2\) and \(|y'|\le C_n\). On either
sign of \(s-t_0\), the map
\[
 (y',s)\longmapsto (Y',Y_n)
           =\big(y'/(s-t_0),\,1/(s-t_0)\big)
\]
and its inverse are uniformly Lipschitz there. It sends the line
\(y'=x'+(s-t_0)v\) to \(Y'=v+Y_n x'\), so the direction
\((x',1)/|(x',1)|\) is independent of \(v\). The image tube is
covered by a bounded number of unit tubes of width \(C_n\rho\) in
that direction. For each sign, transport \(f\) restricted to
\(5Q\cap\{c_n/4\le\pm(s-t_0)\le2\}\) by this map and extend by
zero; this is a single input for all \(x'\) in the slice. Its
\(L^n\) norm is at most \(C\|f\|_{L^n(5Q)}\) by the Jacobian
bounds. The map from \(x'\) to directions has area Jacobian bounded
below. Applying the Kakeya estimate in each slice, integrating in
\(t_0\), then translating and summing over
cubes, gives \(\|D_\rho f\|_n\lesssim L(1/\rho)^{2A}\|f\|_n\).
For \(\rho\) bounded below the same estimate follows by averaging
over a fixed ball about \(x\).

Return to an arbitrary segment through \(x\). Separate the part
within \(O(\delta)\) of \(x\), then divide each of the two rays into
distances \(s/2\le u\le s\) for dyadic \(\delta\le s\le1\).
Relative to the full tube volume, the innermost ball has weight
\(O(\delta)\), and a piece at distance \(s\) has weight \(O(s)\).
By dilation the latter's maximal average has the norm of
\(D_{\delta/s}\). Thus the operator norm is bounded by
\(\displaystyle C\delta+C\sum_s s L(s/\delta)^{2A}
\lesssim L(1/\delta)^{2A}\), proving \eqref{eq:kakeya-nikodym}.
The transverse disk proof above now applies to any Borel set meeting
a suitable unit segment through each point of a positive measure
set in length at least \(\lambda>0\). Such a set has Hausdorff
dimension \(n\). The Nikodym maximal estimate and this dimension conclusion
hold locally on a manifold of
constant sectional curvature for geodesic segments of fixed short
length in a compact subset of a straightening chart: there lengths,
tube widths and volumes are comparable before and after straightening
\cite[Proposition 2.3]{GaoLiuXi2025}.

For a different application, fix \(n\ge3\) and a real \(d\times d\)
matrix \(B\). Consider compact sets containing for every
\(y\in[-1,1]^d\) one curve
\[
           \{(w(y)+(tI_d+t^2 B)y,t):|t|\le\tau_B\},
\]
where \(\tau_B>0\) is fixed sufficiently small and \(w(y)\) belongs
to the fixed sufficiently large compact ball of
\cite[Equation~(1.16)]{Nadjimzadah2026}. Nadjimzadah's lifting
theorem \cite[Theorem 1.3]{Nadjimzadah2026} now applies, since
Corollary~\ref{cor:quant-kakeya} gives full Kakeya dimension in
\(n+\rank(B^2)\) dimensions. The minimum Hausdorff dimension in
this curve family is attained and is
\[
 n-\rank B+\rank(B^2)=n-\dim(\operatorname{im}B\cap\ker B);
\]
the last equality is rank--nullity for \(B\) on its image.
